\documentclass[../main.tex]{subfiles}
\begin{document}
\chapter{Dual Spaces}
\section{Definition}
\begin{definition}[Dual Space]
  For an $\F$-vector space $V$, its \textit{dual space} is defined to be:
  \[
    V^{*} = \mathcal{L}(V, \F) = \{(\theta: V \to \F): \theta \text{ is linear}\}
  \]
  which is in itself another $\F$-vector space.
\end{definition}
\begin{example}
  \begin{enumerate}
    \item For $V = \R^{4}$, the elements $\theta_1, \theta_2, \theta_3 \in V^{*}$ defined as:
      \begin{align*}
        \theta_1(v) &= 2(v_2 + v_4) - v_3  \\
        \theta_2(v) &= 2(v_2 - v_4) - v_3  \\
        \theta_3(v) &= 2v_2 - 4v_3 - 3v_4  \\
      \end{align*}
      are linearly dependent.
    \item $\tr \in (\mathcal{M}_{n \times n}(\F))^{*}$
    \item Consider the $\F$-vector space $V = C([0, 1])$, the map:
      \[
        \theta(f) = \int_{0}^{1} f(t)e^{-t} \d{t}
      \]
      has $\theta \in V^{*}$
  \end{enumerate}
\end{example}
\begin{lemma}
  If $V$ is a finite dimensional $\F$-vector space, then $V^{*}$ is also finite dimensional with $\dim(V^{*}) = \dim(V)$. \label{dualSize}
\end{lemma}
\begin{proof}
  By \cref{linMapDim}, $\dim(V^{*}) = \dim(\mathcal{L}(V, \F)) = \dim(V) \dim(\F) = \dim(V)$.
\end{proof}
\section{Dual Basis}
\begin{definition}[Dual Basis]
  Let $V$ be an $\F$-vector space and let $B$ be a basis for $V$.
  For each $b \in B$, define $b^{*} \in V^{*}$ by:
  \[
    b^{*}\left(\finitesum_{c \in B} \lambda_c c\right) = \lambda_b
  \]
  i.e. this the unique linear extension (\cref{linMapBasis}) of $b^{*}(c) = \delta_{b c}$ for $c \in B$.
  We then \textit{dual basis} of $B$ as:
  \[
    B^{*} = \{b^{*}: b \in B\} \subseteq V^{*}
  \]
\end{definition}
\begin{proposition}[Properties of the dual basis]
  If $B$ is a basis for an $\F$-vector space $V$, then its dual basis $B^{*}$ satisfies: \label{dualBasis}
  \begin{enumerate}
    \item $B^{*}$ is linearly independent
    \item If $V$ is finite dimensional then $B^{*}$ is a basis for $V^{*}$.
  \end{enumerate}
\end{proposition}
\begin{proof}[\textbf{i}]
  Suppose that we have a linear combination of elements of $B^{*}$ equal to $\vec{0}$:
  \[
    \finitesum_{b \in B} \lambda_b b^{*} = \vec{0}
  \]
  This is a function in $V^{*} = \mathcal{L}(V, \F)$, passing in any $c \in B$ yields:
  \[
    0 = \finitesum \lambda_b b^{*}(c) = \finitesum \lambda_b \delta_{b c} = \lambda_c
  \]
  so all the coefficients are 0, i.e. $B^{*}$ is linearly independent.
\end{proof}
\begin{proof}[\textbf{ii} - Using dimensions]
  By \cref{dualBasis}, $\dim(V^{*}) = \dim(V) = |B|$.
  From \textbf{(i)}, we know $B^{*}$ is linearly independent, so by \cref{dimIneq}, $B^{*}$ is a basis for $V^{*}$.
\end{proof}
\begin{proof}[\textbf{ii} - Constructing the span]
  We can show that any $\theta \in V^{*}$ lies in the span of $B^{*}$.

  Suppose $\theta \in V^{*}$ and define the coefficients $\lambda_b = \theta(b) \in \F$ for $b \in B$.
  Define $\overline{\theta} \in V^{*}$ as:
  \[
    \overline{\theta} = \finitesum_{b \in B} \lambda_b b^{*} \in \langle B^{*} \rangle
  \]
  so that for $c \in B$:
  \[
    \overline{\theta}(c) = \finitesum_{b \in B} \lambda_b b^{*}(c) = \lambda_c = \theta(c)
  \]
  Thus $\overline{\theta}$ and $\theta$ agree on the basis $B^{*}$ so by \cref{linMapBasis} $\overline{\theta} = \theta \in \langle B^{*} \rangle$.
  Thus since this is true for all $\theta$, $\langle B^{*} \rangle = V^{*}$.
\end{proof}
\begin{remark}
  If $V$ is infinite dimensional, then the second proof of \textbf{(ii)} above could run into issues if infinitely many $\lambda_b$ are non-zero as then we could not construct $\overline{\theta}$ as a finite linear combination.
\end{remark}
\begin{corollary}
  If $V$ is a finite dimensional $\F$-vector space, then $V \cong V^{*}$.
\end{corollary}
\begin{proof}
  From \cref{dualSize}, $\dim(V) = \dim(V^{*})$ and so by \cref{finiteDimIso}, we have $V \cong V^{*}$.
\end{proof}
\begin{example}
  Consider the case of $V = P(\R)$, which is infinite dimensional.
  We will show on Example Sheet 2 that $V^{*} \cong \R^\N$ but we also showed on Example Sheet 1 that $P(\R) \centernot\cong \R^\N$ and so the above does not always hold if $V$ is infinite dimensional.
\end{example}
If $V$ is finite dimensional and $B = \{b_i: 1 \leq i \leq n\}$, then we can use the basis and dual basis to express the application of $\theta \in V^{*}$ as a dot product.
Since $B$ is a basis for $V$ and $B^{*} = \{b^{*}_i: 1 \leq i \leq n\}$ is a basis for $V^{*}$, we can write:
\[
  v = \sum_{i = 1}^{n} \lambda_i b_i \text{ and } \theta = \sum_{j = 1}^{n} \mu_i b^{*}_i
\]
Thus:
\begin{align*}
  \theta(v) &= \sum_{j = 1}^{n} \mu_i b^{*}_i\left[\sum_{i = 1}^{n} \lambda_i b_i\right] \\
            &= \sum_{i, j = 1}^{n} \lambda_i \mu_j b^{*}_j(b_i) \\
            &= \sum_{i = 1}^{n} \lambda_i \mu_i
\end{align*}
From the linear combinations we see that:
\[
  \begin{pmatrix}
  \lambda_1 \\
  \vdots \\
  \lambda_n \\
  \end{pmatrix} = [v]_B \text{ and } \begin{pmatrix}
  \mu_1 \\
  \vdots \\
  \mu_n \\
  \end{pmatrix} = [\theta]_{B^{*}}
\]
so $\theta(v) = ([\theta]_{B^{*}})^{\trans} [v]_{B}$.
Note that for this to work, we need $B^{*}$ to be the dual basis for $B$ specifically and not just any basis for $V^{*}$
\section{Annihilators}
\begin{definition}[Annihilator]
  For a subset $S \subseteq V$, we define the \textit{annihilator} of $S$ in $V^{*}$ as:
  \[
    S^{0} = \{\theta \in V^{*}: \theta(s) = 0\ \forall s \in S\}
  \]
\end{definition}
\begin{lemma}[Properties of annahilators]
  For any subsets $S, T \subseteq V$, we have: \label{annaProp}
  \begin{enumerate}
    \item $S^{0} \leq V^{*}$
    \item If $S \subseteq T$, then $T^{0} \leq S^{0}$, i.e. taking annihilators is \textit{inclusion reversing}.
    \item $S^{0} = \langle S \rangle^{0}$
    \item $V^{0} = \{\vec{0}_{V^{*}}\}$ and $\{\vec{0}_{V}\}^{0} = V^{*}$.
  \end{enumerate}
\end{lemma}
\begin{proof}[\textbf{i}]
  Suppose $\theta_1, \theta_2 \in S^{0}$ and $\lambda \in \F$, then for any $s \in S$:
  \[
    (\theta_1 + \lambda \theta_2)(s) = \theta_1(s) + \lambda \theta_2(s) = 0
  \]
  so $\theta_1 + \lambda \theta_2 \in S^{0}$.
  Thus by the subspace test (\cref{subspaceTest}), $S^{0} \leq V^{*}$.
\end{proof}
\begin{proof}[\textbf{ii}]
  By \textbf{(i)}, it is sufficient to check that $T^{0} \subseteq S^{0}$, which is immediate from the definition.
\end{proof}
\begin{proof}[\textbf{iii}]
  We have $S \subseteq \langle S \rangle$, so by \textbf{(ii)}, $\langle S \rangle^{0} \subseteq S^{0}$.

  Conversely, take any $\theta \in S^{0}$ and any $v \in \langle S \rangle$.
  We can write:
  \[
    v = \finitesum_{s \in S} \lambda_s s \in \langle S \rangle
  \]
  and so:
  \[
    \theta(v) = \finitesum_{s \in S} \lambda_s \theta(s) = \finitesum_{s \in S} \lambda_s \cdot 0 = 0
  \]
  i.e. $\theta \in \langle S \rangle^{0}$.
  So $S^{0} \subseteq \langle S \rangle^{0}$, thus $S^{0} = \langle S \rangle^{0}$.
\end{proof}
\begin{proof}[\textbf{iv}]
  \begin{itemize}
    \item If $\theta \in V^{0}$, then $\theta(v) = 0$ for all $v \in V$, so $\theta = \vec{0}_{V^{*}}$.
    \item If $\theta \in V^{*}$, then $\theta(\vec{0}_{V}) = 0$ since $\theta$ is linear, so $\theta \in \{\vec{0}_{V}\}^{0}$, i.e. $V^{*} \subseteq \{\vec{0}_{V}\}^{0}$.
      By definition $\{\vec{0}_{V}\}^{0} \subseteq V^{*}$ and so $V^{*} = \{\vec{0}_{V}\}^{0}$
  \end{itemize}
\end{proof}
\begin{proposition}[Dimensions of annahilators]
  If $V$ is a finite dimensional $\F$-vector space with a subspace $U \leq V$, then: \label{annaDim}
  \[
    \dim(V) = \dim(U) + \dim(U^{0})
  \]
\end{proposition}
\begin{proof}
  Let $B_U = \{v_1, \cdots, v_k\}$ be a basis for $V$ which we extend using \cref{basisExtension} to a basis $B_V = \{v_1, \ldots, v_n\}$ for $V$, leading to a dual basis $B^{*}_V = \{v^{*}_1, \ldots, v^{*}_n\}$.
  It is sufficient to prove that $\widetilde{B} = \{v_{k + 1}^{*}, \ldots, v^{*}_n\}$  is a basis for $U^{0}$ as this has size $\dim(V) - \dim(U)$.

  \textbf{Subset}\par
  We first need to check that each of the vectors in $\widetilde{B}$ is in $U^{0}$.
  For $i \leq k$ and $j \geq k + 1$, we have $v^{*}_j(v_i) = 0$ as we cannot have $i = j$, thus $v^{*}_j$ eliminates $B_U$ so $v^{*}_j \in (B_{U})^{0} = \langle B_U \rangle^{0}$ by \cref{annaProp}.
  But $B_U$ is a basis for $U$ so $v^{*}_j \in \langle B_U \rangle^{0} = U^{0}$, i.e. $\widetilde{B} \subseteq U^{0}$.

  \textbf{Linearly independent}\par
  $\widetilde{B}$ is a subset of $B^{*}_V$ which is a basis so $\widetilde{B}$ is linearly independent.

  \textbf{Spanning}\par
  Suppose $\theta \in U^{0} \leq V^{*}$.
  By \cref{dualBasis}, $B^{*}$ is a basis for $V^{*}$ and so we can write:
  \[
    \theta = \sum_{j = 1}^{n} \lambda_j v^{*}_j
  \]
  For $i \leq k$ we have $v_i \in U$ so since $\theta \in U^{0}$, we have $\theta(v_i) = 0$ and thus:
  \[
    0 = \theta(v_i) = \sum_{j = 1}^{n} \lambda_j v^{*}_j(v_i) = \sum_{j = 1}^{n} \lambda_j \delta_{i j} = \lambda_i
  \]
  so $\lambda_i = 0$ for $i \leq k$.
  Hence $\theta = \sum_{j = k + 1}^{n} \lambda_j v^{*}_j$, which is a linear combination of only vectors in $\widetilde{B}$, so $\widetilde{B}$ spans $U^{0}$.
\end{proof}
\begin{remark}
  By how we proved the above, if $U \leq V$, then there exists $X \leq V^{*}$ such that:
  \[
    U = \{v \in V: \forall \theta \in X,\ \theta(v) = 0\}
  \]
  with $\dim(X) = \dim(V) - \dim(U)$.
\end{remark}
\begin{proposition}[Annahilators of subspace sums]
  If $V$ is an $\F$-vector space and $U, W \leq V$, then:
  \begin{enumerate}
    \item $(U + W)^{0} = U^{0} \cap W^{0}$
    \item $U^{0} + W^{0} \leq (U \cap W)^{0}$
    \item If $V$ is finite dimensional, then we have equality in \textbf{(ii)}
  \end{enumerate}
\end{proposition}
\begin{proof}[\textbf{i}]
  For any $\theta \in V^{*}$:
  \begin{align*}
    \theta \in (U + W)^{0} &\iff \theta(u + w) = 0\ \forall u \in U, w \in W \\
                           &\iff \theta(u) + \theta(w) = 0\ \forall u \in U, w \in W \\
                           &\iff \theta(u) = \theta(w) = 0\ \forall u \in U, w \in W \text{ by taking either $u$ or $w$ to be $\vec{0}$} \\
                           &\iff \theta \in U^{0} \cap W^{0}
  \end{align*}
  So $(U + W)^{0} = U^{0} \cap W^{0}$.
\end{proof}
\begin{proof}[\textbf{ii}]
  We see that $U \cap W \leq U, W$, so by \cref{annaProp} \textbf{(ii)} taking annihilators is inclusion reversing, thus $U^{0}, W^{0} \leq (U \cap W)^{0}$.
  Therefore by closure, $U^{0} + W^{0} \leq (U \cap W)^{0}$, as required.
\end{proof}
\begin{proof}[\textbf{iii}]
  So $\dim(V) = n$, thus using the sum-intersection formula (\cref{sumInter}), we have:
  \begin{align*}
    \dim(U^{0} + W^{0}) &= \dim(U^{0}) + \dim(W^{0}) - \dim(U^{0} \cap W^{0}) \\
                        &= \dim(U^{0}) + \dim(W^{0}) - \dim((U + W)^{0}) \text{ by \textbf{(i)}} \\
                        &= n - \dim(U) - \dim(W) + \dim(U + W) \text{ by \cref{annaDim}} \\
                        &= n - \dim(U \cap W) \text{ by \cref{sumInter}} \\
                        &= \dim((U \cap W)^{0}) \text{ by \cref{annaDim}}
  \end{align*}
  From \textbf{(ii)}, we have $U^{0} + W^{0} \leq (U \cap W)^{0}$.
  Thus using \cref{finiteSubspace}, $U^{0} + W^{0} = (U \cap W)^{0}$.
\end{proof}
\section{Dual Maps}
\begin{definition}[Dual map]
  For $\F$-vector spaces $V, W$ and a linear map $\alpha \in \mathcal{L}(V, W)$, we define the \textit{dual map} of $\alpha$ as:
  \[
    \alpha^{*}: W^{*} \to V^{*},\ \theta \mapsto \theta \circ \alpha
  \]
\end{definition}
\begin{lemma}[Properties of dual maps]
  If $U, V, W$ are $\F$-vector spaces, $\alpha, \beta \in \mathcal{L}(V, W)$, $\gamma \in \mathcal{L}(U, V)$, and $\lambda \in \F$ then:
  \begin{enumerate}
    \item $\alpha^{*} \in \mathcal{L}(W^{*}, V^{*})$
    \item $(\alpha + \lambda \beta)^{*} = \alpha^{*} + \lambda \beta^{*}$
    \item $(\alpha \circ \gamma)^{*} = \gamma^{*} \circ \alpha^{*}$
    \item If $\beta$ is an isomorphism, then so is $\beta^{*}$ with $(\beta^{*})^{-1} = (\beta^{-1})^{*}$
  \end{enumerate}
\end{lemma}
\begin{remark}
  The above properties are similar to the those of the transpose of a matrix, we will see later why this is the case.
\end{remark}
\begin{remark}[Advice]
  When trying to show equality of maps it is usually easiest to show that the maps agree when applied to an arbitrary element of the domain.
\end{remark}
\begin{proof}[\textbf{i}]
  Let $\theta, \eta \in W^{*}$ and $\mu \in \F$ be arbitrary.
  For any $v \in V$ we have:
  \begin{align*}
    \alpha^{*}(\theta + \mu \eta)[v] &= (\theta + \mu \eta)(\alpha(v)) \\
                                     &= \theta(\alpha(v)) + \mu \eta(\alpha(v)) \\
                                     &= \alpha^{*}(\theta)[v] + \mu \alpha^{*}(\eta)[v] \\
                                     &= (\alpha^{*}(\theta) + \mu \alpha^{*}(\eta))[v]
  \end{align*}
  Since this is true for any $v \in V$, $\alpha^{*}(\theta + \mu \eta) = \alpha^{*}(\theta) + \mu \alpha^{*}(\eta)$, i.e. $\alpha^{*}$ is linear.
\end{proof}
\begin{proof}[\textbf{ii}]
  For $\theta \in W^{*}$ and $v \in V$ we have:
  \begin{align*}
    (\alpha + \lambda \beta)^{*}(\theta)[v] &= \theta((\alpha + \lambda \beta)(v)) \\
                                            &= \theta(\alpha(v) + \lambda \beta(v)) \\
                                            &= \theta(\alpha(v)) + \lambda \theta(\beta(v)) \\
                                            &= \alpha^{*}(\theta)[v] + \lambda \beta^{*}(\theta)[v] \\
                                            &= (\alpha^{*} + \lambda \beta^{*})(\theta)[v]
  \end{align*}
  This is true for any $v \in V$ so $(\alpha + \lambda \beta)^{*}(\theta) = (\alpha^{*} + \lambda \beta^{*})(\theta)$ for any $\theta \in W^{*}$, and so $(\alpha + \lambda \beta)^{*} = \alpha^{*} + \lambda \beta^{*}$.
\end{proof}
\begin{proof}[\textbf{iii}]
  For any $\theta \in W^{*}$:
  \begin{align*}
    (\alpha \circ \gamma)^{*}(\theta) &= \theta \circ (\alpha \circ \gamma) \\
                                      &= (\theta \circ \alpha) \circ \gamma \\
                                      &= (\alpha^{*}(\theta)) \circ \gamma \\
                                      &= \gamma^{*}(\alpha^{*}(\theta)) \\
                                      &= (\gamma^{*} \circ \alpha^{*})(\theta)
  \end{align*}
  Since this is true for any $\theta \in W^{*}$, $(\alpha \circ \gamma)^{*} = \gamma^{*} \circ \alpha^{*}$.
\end{proof}
\begin{proof}[\textbf{iv}]
  For all $\theta \in V^{*}$:
  \[
    (\id_V)^{*}(\theta) = \theta \circ \id_V = \theta = \id_{V^{*}}(\theta)
  \]
  Since this is true for all $\theta \in V^{*}$, $(\id_V)^{*} = \id_{V^{*}}$.

  Hence using \textbf{(iii)}, we have:
  \begin{align*}
    \id_{V^{*}} &= (\beta^{-1} \circ \beta)^{*} = \beta^{*} \circ (\beta^{-1})^{*} \\
    \text{ and } \id_{V^{*}} &= (\beta \circ \beta^{-1})^{*} = (\beta^{-1})^{*} \circ \beta^{*}
  \end{align*}
  and so $(\beta^{-1})^{*} = (\beta^{*})^{-1}$.
\end{proof}
\begin{proposition}
  If $V$ and $W$ are finite dimensional with bases $B$ and $C$ respectively and $\alpha \in \mathcal{L}(V, W)$, then: \label{starTrans}
  \[
    [\alpha^{*}]^{C^{*}}_{B^{*}} = ([\alpha]^{B}_{C})^{\trans}
  \]
\end{proposition}
\begin{proof}
  Let $n = \dim(V)$ and $m = \dim(W)$, i.e. $B = \{b_1, \ldots, b_n\}$ and $C = \{c_1, \ldots, c_m\}$.
  Define the matrix $A = [\alpha]^{B}_{C}$.
  By definition:
  \[
    \alpha(b_i) = \sum_{k = 1}^{m} a_{k i}c_k
  \]
  Similarly, define $A' = [\alpha^{*}]^{C^{*}}_{B^{*}}$, so that:
  \[
    \alpha^{*}(c^{*}_j) = \sum_{\ell = 1}^{n} a_{\ell j}' b^{*}_{\ell}
  \]
  By definition $\alpha^{*}: W^{*} \to V^{*}$ so $\alpha^{*}(c^{*}_j) \in V^{*}$, thus we can pass in $b_i$:
  \begin{align*}
    \alpha^{*}(c^{*}_j)[b_i] &= \sum_{\ell = 1}^{n} a'_{\ell j} b^{*}_{\ell}(b_i) \\
                             &= \sum_{\ell = 1}^{n} a
                             _{\ell j} \delta_{i \ell} \\
                             &= a'_{i j}
  \end{align*}
  Instead, if we use the definition of $\alpha^{*}$:
  \begin{align*}
    \alpha^{*}(c^{*}_j)[b_i] &= (c^{*}_j \circ \alpha)[b_i] \\
                             &= c^{*}_j(\alpha(b_i)) \\
                             &= \sum_{k = 1}^{m} a_{k i} c^{*}_j(c_k) \\
                             &= \sum_{k = 1}^{m} a_{k i} \delta_{j k} \\
                             &= a_{j i}
  \end{align*}
  So $a'_{i j} = a_{j i}$ for all $i, j$, i.e. $[\alpha^{*}]^{C^{*}}_{B^{*}} = ([\alpha]^{B}_{C})^{\trans}$.
\end{proof}
\begin{example}[Change of basis]
  Let $B$ and $C$ be bases for a finite $\F$-vector space $V$.
  Consider the change of basis matrix (\cref{cobFormula}) from $B$ to $C$, $P = [\id_V]^{B}_{C} \in \GL_n(\F)$.
  To find the change of basis matrix from $B^{*}$ to $C^{*}$, $[\id_{V^{*}}]^{B^{*}}_{C^{*}}$:
  \[
    P^{\trans} = [\id_{V^{*}}]^{C^{*}}_{B^{*}} \implies (P^{\trans})^{-1} =  [\id_{V^{*}}]^{B^{*}}_{C^{*}}
  \]
  Thus we need to take both the transpose and the inverse to find the corresponding change of basis matrix from $B^{*}$ to $C^{*}$.
\end{example}
\begin{proposition}[Annahilators of kernels and images]
  If $V$ and $W$ are (not necessarily finite dimensional) $\F$-vector spaces with $\alpha \in \mathcal{L}(V, W)$, then:
  \begin{enumerate}
    \item $\ker(\alpha^{*}) = (\im \alpha)^{0}$
    \item $\im(\alpha^{*}) \leq (\ker \alpha)^{0}$
    \item If $V$ and $W$ are finite dimensional, then $\rank(\alpha^{*}) = \rank(\alpha)$ and we have equality in \textbf{(ii)}
  \end{enumerate}
\end{proposition}
\begin{proof}[\textbf{i}]
  For $\theta \in W^{*}$:
  \begin{align*}
    \theta \in \ker(\alpha^{*}) &\iff \theta \circ \alpha = \vec{0}_{V^{*}} \\
                                &\iff \theta(\alpha(v)) = 0\ \forall v \in V \\
                                &\iff \theta \in (\im \alpha)^{0}
  \end{align*}
  Thus $\ker(\alpha^{*}) = (\im \alpha)^{0}$.
\end{proof}
\begin{proof}[\textbf{ii}]
  Suppose $\eta \in \im(\alpha^{*})$, so there exists $\theta \in W^{*}$ such that $\eta = \alpha^{*}(\theta)$.

  For any $v \in \ker \alpha$:
  \begin{align*}
    \eta(v) &= \alpha^{*}(\theta)[v] \\
            &= \theta(\alpha(v)) \\
            &= \theta(0) \\
            &= 0
  \end{align*}
  So $\eta \in (\ker \alpha)^{0}$.
  Thus $\im(\alpha^{*}) \subseteq (\ker \alpha)^{0}$, but $\im(\alpha^{*}), (\ker \alpha)^{0} \leq V^{*}$ so $\im(\alpha^{*}) \leq (\ker \alpha)^{0}$.
\end{proof}
\begin{proof}[\textbf{iii}]
  Since $V$ and $W$ are finite, we can always pick basis $B$ and $C$ for $V$ and $W$ respectively.
  Using \cref{starTrans}, we have:
  \[
    A = [\alpha]^{B}_{C} \text{ and } A^{\trans} = [\alpha^{*}]^{C^{*}}_{B^{*}}
  \]
  Then, from \cref{rowColRank}, the row rank of a matrix is equal to its column rank and so:
  \[
    \rank \alpha = \colrank(A) = \rowrank(A) = \colrank(A^{\trans}) = \rank(\alpha^{*})
  \]

  We then use \cref{annaDim}:
  \begin{align*}
    \dim((\ker \alpha)^{0}) &= \dim(V) - \nullity(\alpha) \text{ by \cref{annaDim}} \\
                            &= \rank \alpha \text{ by \cref{RNT}} \\
                            &= \dim(\im \alpha)
  \end{align*}
  From \textbf{(ii)}, $\im(\alpha^{*}) \leq (\ker \alpha)^{0}$ so by \cref{finiteSubspace}, $\im(\alpha^{*}) = (\ker \alpha)^{0}$.
\end{proof}
We can also prove \textbf{(iii)} another way without needing \cref{starTrans}.
\begin{proof}[\textbf{iii} -- Alternative]
  Suppose that $\rank \alpha^{*} < \rank \alpha$, then
  \begin{align*}
    \dim(W) &= \dim(W^{*}) \text{ by \cref{dualSize}} \\
            &= \rank(\alpha^{*}) + \nullity(\alpha^{*}) \text{ by \cref{RNT}} \\
            &< \rank \alpha + \nullity(\alpha^{*}) \text{ by supposition} \\
            &= \rank \alpha + \dim(\im(\alpha)^{0}) \text{ by \textbf{(i)}} \\
            &= (\dim(V) - \nullity \alpha) + (\dim(W) - \rank \alpha) \text{ by \cref{annaDim}} \\
            &= \dim(W) \text{ by \cref{RNT}}
  \end{align*}
  which is a contradiction, so $\rank \alpha^{*} \geq \rank \alpha$.

  As in the first proof, we have that $\dim((\ker \alpha)^{0}) = \rank \alpha$, and so using \textbf{(ii)}, we conclude that $\im(\alpha^{*}) = (\ker \alpha)^{0}$.
\end{proof}
\begin{remark}
  Since this proof does not utilise $\rowrank A = \colrank A$ (\cref{rowColRank}), we can use the above to provide an alternative proof for this equality.

  For any $A \in \mathcal{M}_{m \times n}(\F)$, we can construct the map:
  \[
    \alpha: \F^n \to \F^m,\ v \mapsto Av
  \]
  so that $A = [\alpha]^{\varepsilon_n}_{\varepsilon_m}$ where $\varepsilon_n$ and $\varepsilon_m$ are the standard bases for $\F^n$ and $\F^m$ respectively.
  Thus by \cref{starTrans} $A^{\trans} = [\alpha^{*}]^{E_m^{*}}_{E_n^{*}}$ and so using \textbf{(iii)}:
  \[
    \colrank(A) = \rank \alpha = \rank(\alpha^{*}) = \colrank(A^{\trans}) = \rowrank(A)
  \]
\end{remark}
\section{Double Dual}
\begin{definition}[Double dual]
  The \textit{double dual} of an $\F$-vector space $V$ is the dual of the dual space, that is:
  \[
    V^{**} = (V^{*})^{*} = \mathcal{L}(V^{*}, \F)
  \]
\end{definition}
\begin{remark}
  If $V$ is finite dimensional, then by \cref{dualSize}, $\dim(V^{*}) = \dim(V)$ and $\dim(V^{*}) = \dim(V^{**})$, so $\dim(V) = \dim(V^{**})$.
  Thus by \cref{finiteDimIso}, $V \cong V^{**}$.
\end{remark}
\begin{theorem}[Evaluation map]
  If $V$ is an $\F$-vector space, the \textit{evaluation map} or \textit{natural isomorphism} $V \to V^{**}$ is defined as: \label{evalMap}
  \[
    \varepsilon: V \to V^{**},\ v \mapsto (\varepsilon(v): V^{*} \to \F,\ \theta \mapsto \theta(v)) \in \mathcal{L}(V^{*}, \F)
  \]
  is a linear map.
  Moreover, if $V$ is finite dimensional then $\varepsilon$ is an isomorphism.
\end{theorem}
\begin{remark}
  More compactly, $\varepsilon$ satisfies $\varepsilon(v)[\theta] = \theta(v)$ for all $v \in V$ and $\theta \in V^{*}$.
\end{remark}
\begin{proof}\par
  \textbf{Linearity}\par
  Let $v, w \in V$ and $\lambda \in \F$.
  For any $\theta \in V^{*}$, we have:
  \begin{align*}
    \varepsilon(v + \lambda w)[\theta] &= \theta(v + \lambda w) \\
                                       &= \theta(v) + \lambda \theta(w) \text{ as $\theta$ is linear} \\
                                       &= \varepsilon(v)[\theta] + \lambda \varepsilon(w)[\theta] \\
                                       &= (\varepsilon(v) + \lambda\varepsilon(w))[\theta]
  \end{align*}
  This holds for all $\theta \in V^{*}$ and so $\varepsilon(v + \lambda w) = \varepsilon(v) + \lambda \varepsilon(w)$, i.e. $\varepsilon$ is linear.

  Suppose further that $V$ is finite dimensional.\par
  \textbf{Injective}\par
  Suppose that $\ker(\varepsilon) \neq \{\vec{0}_{V}\}$, thus there exists $v \in \ker(\varepsilon) \setminus \{\vec{0}_{V}\}$.
  So for all $\theta \in V^{*}$:
  \[
    \varepsilon(v)[\theta] = \theta(v) = 0
  \]
  For contradiction, we want to find a $\theta \in V^{*}$ for which this does not hold.

  We use \cref{basisExtension} to extend $\{v\}$ to a basis $B = \{v = v_1, v_2, \ldots, v_n\}$ for $V$, and let $B^{*} = \{v^{*}_1, \ldots, v^{*}_n\}$ be the dual basis to $B$.
  Taking $\theta= v^{*}_1$, we see that $1 = \theta(v) = 0$, which is a contradiction, so $\ker(\varepsilon) = \{\vec{0}\}$, i.e. $\varepsilon$ is injective.

  \textbf{Surjectivity}\par
  Since $\dim(V^{**}) = \dim(V)$, by \cref{linPidge}, we have that $\varepsilon$ is an isomorphism.
\end{proof}
\begin{remark}
  We knew already by a dimension argument that if $V$ is finite dimensional, then $V \cong V^{**}$, however the isomorphisms provided by this argument will always depend on a choice of basis.

  This makes $\varepsilon$ a particularly useful isomorphism to work with as it does not depend on any choice of basis.
\end{remark}
\begin{proposition}
  If $V$ is finite dimensional with $\dim(V) = n$, then every basis $C$ of $V^{*}$ is dual to some basis for $V$.
\end{proposition}
\begin{proof}
  Write $C$ as $C = \{\theta_1, \ldots, \theta_n\} \subseteq V^{*}$.
  We can find the dual basis $C^{*} = \{\theta^{*}_1, \ldots, \theta^{*}_n\}$ which is a basis for $V^{**}$.

  Define the set $B = \{v_1, \ldots, v_n\}$ by $v_i = \varepsilon^{-1}(\theta^{*}_i)$.
  By \cref{evalMap}, $\varepsilon$ is an isomorphism and so since $C^{*}$ is a basis for $V^{**}$, we must have that $B = \{v_1, \ldots, v_n\}$ is a basis for $V$.

  We now need to check that $B^{*} = C$.
  Consider the following:
  \begin{align*}
    \theta_i(v_j) &= \varepsilon(v_j)[\theta] \text{ by definition of $\varepsilon$} \\
                  &= \theta^{*}_j(\theta_i) \text{ by definition of $v_j$} \\
                  &= \delta_{i j} \text{ by definition of the dual basis}
  \end{align*}
  So for each $i$, $\theta_i$ agrees with $v^{*}_i$ on each element $v_j \in B$, thus by \cref{linMapBasis}, $\theta_i = v^{*}_i$.
  Therefore, $B^{*} = \{v^{*}_1, \ldots, v^{*}_n\} = \{\theta_1, \ldots, \theta_n\} = C$.
\end{proof}
\begin{proposition}
  If $V$ is finite dimensional with $\dim(V) = n$, then for any $U \leq V$, $\varepsilon(U) = (U^{0})^{0}$
\end{proposition}
\begin{remark}
  $\varepsilon(U)$ is the image of the subspace $U$ under $\varepsilon$ and so is a subset of $V^{**}$ along with $(U^{0})^{0}$. \cref{doubleAnna}
\end{remark}
\begin{proof}
  If $u \in U$, then for all $\theta \in U^{0}$:
  \[
    \varepsilon(u)[\theta] = \theta(u) = 0
  \]
  i.e. $\varepsilon(u) \in V^{*}$ annihilates all $\theta \in U^{0}$, so $\varepsilon(u) \in (U^{0})^{0}$ for all $u \in U$.
  Thus $\varepsilon(U) \leq (U^{0})^{0}$.

  \begin{align*}
    \dim((U^{0})^{0}) &= \dim(V^{*}) - \dim(U^{0}) \text{ by \cref{annaDim}} \\
                      &= \dim(V^{*}) - (\dim(V) - \dim(U)) \text{ by \cref{annaDim}} \\
                      &= \dim(U) \text{ by \cref{dualSize}} \\
                      &= \dim(\varepsilon(U)) \text{ as $\varepsilon$ is an isomorphism}
  \end{align*}
  So by \cref{finiteSubspace}, $\varepsilon(U) = (U^{0})^{0}$.
\end{proof}
\begin{remark}
  For any subspace $X \leq V^{*}$, we have $X^{0} \leq V^{**}$.
  This subspace $X^{0}$ will have a preimage under $\varepsilon$, namely $X_0 = \varepsilon^{-1}(X^{0}) \leq V$.
  We can find a more explicit expression for $X_0$:
  \begin{align*}
    X_0 &= \{v \in V: \varepsilon(v) \in X^{0}\} \\
        &= \{v \in V: \varepsilon(v)[\theta] = 0\ \forall \theta \in X\} \\
        &= \{v \in V: \theta(v) = 0\ \forall \theta \in X\} \\
        &= \{v \in V: v \in \ker \theta\ \forall \theta \in X\} \\
        &= \bigcap_{\theta \in X} \ker \theta
  \end{align*}
  Since $\varepsilon$ is an isomorphism, we further have that $\dim(X_0) = \dim(X^{0}) = n - \dim(X)$.

  In particular, given a subspace $U \leq V$, if we let $X = U^{0} \leq V^{*}$, then
  \begin{align*}
    X_0 &= \varepsilon^{-1}((U^{0})^{0}) \\
        &= \varepsilon^{-1}(\varepsilon(U)) \text{ by \cref{doubleAnna}} \\
        &= U \text{ as $\varepsilon$ is an isomorphism}
  \end{align*}
  So for any subspace $U \leq V$, $U$ is exactly the zero locus of its annihilator $U^{0} \leq V^{*}$, i.e.:
  \[
    U = \bigcap_{\theta \in U^{0}} \ker \theta
  \]
  A priori, this may not have been true as there could have been elements of $V \setminus U$ which just happened to be sent to $0$ by all $\theta \in U^{0}$.

  For a concrete example of this, if $V = \F^n$, then we can express any subspace $U \leq V$ as a system of linear equations in the components of the vectors under a choice of basis.
\end{remark}
\end{document}
